% Source: https://artofproblemsolving.com/community/c3h3097711_rigor_psa_bust_that_fakesolve
% Stable identity: aops-topic-3097711-q02
Existence follows by repeated Euclidean division by 10: write $n=10q+r$ with unique $0\le r\le9$, and repeat on the smaller quotient until it is zero. Reading the remainders backwards gives the digits. For uniqueness, the final digit must equal the remainder modulo 10. Subtract that digit and divide by 10; the same argument determines every earlier digit. Termination and the no-leading-zero rule ensure that two representations coincide.
